% ============================================================================
% amplifier-bundle-research — white paper (v2, LaTeX)
% Audience: patent attorneys (Persona A).
% Compiled via pdflatex + bibtex.
% Session: e623279d-44ef-458d-b7b7-ddc8f73e105c
% Figures generated by Graphviz from DOT sources in ./figures/.
% ============================================================================
\documentclass[11pt,letterpaper]{article}

\usepackage[utf8]{inputenc}
\usepackage[T1]{fontenc}
\usepackage[margin=1in]{geometry}
\usepackage{amsmath,amssymb}
\usepackage{graphicx}
\usepackage{booktabs}
\usepackage{longtable}
\usepackage{microtype}
\usepackage[numbers,square]{natbib}
\usepackage{xcolor}
\usepackage{enumitem}
\usepackage{parskip}
\usepackage{tabularx}
\usepackage[hidelinks]{hyperref}
\usepackage[capitalize,noabbrev]{cleveref}
\usepackage{caption}
\captionsetup{labelfont=bf,font=small}

\graphicspath{{figures/}}

% Draft-mode annotations (kept visible in v2; strip for v1.0 release)
\newcommand{\exploratory}[1]{\textcolor{orange}{\textbf{[EXPLORATORY]} #1}}

\title{A Methodology Scaffold for Patent Drafting\\
\large How \texttt{amplifier-bundle-research} v0.4.0 makes an attorney's
own discipline visible, auditable, and portable, and what it does not
yet do}

\author{Michael Jabbour\\
\small{Author of the bundle described herein}}

\date{April 23, 2026 \\ \small{Version 0.2, LaTeX rewrite with dogfood diagrams}}

\begin{document}
\maketitle

% ============================================================================
\section*{Disclosure: first-person capability description}
% ============================================================================

The author of this paper wrote the bundle it describes. The paper is
a first-person capability description, not an independent capability
evaluation. Readers should treat the measured claims as walkthrough
evidence produced by the bundle's author, and weigh them accordingly.
The author benefits reputationally (not financially in v0.4.0; the
bundle is open-source under Apache-2.0) from adoption. A v1.0
milestone has been reserved for real-attorney validation by
practitioners the author does not employ or direct; that validation
has not yet been performed.

% ============================================================================
\section*{Executive Summary}
% ============================================================================

Patent prosecution is an evidence-defensibility craft. An attorney
drafting a \S112-compliant specification and a \S103-resistant claim
set is doing work structurally similar to a scientific author drafting
a methods section: committing to a scope and a standard of evidence
before the examiner (or the reviewer) tests it.

The problem this paper addresses is specific. When attorneys use
general-purpose AI (ChatGPT, Claude, Gemini) to accelerate first-draft
work, the methodology steps the human was silently carrying (search
scope discipline, overclaiming detection, enablement self-check) get
quietly skipped. The draft reads well. The examiner finds the gap
later, when amendment space is expensive.

\texttt{amplifier-bundle-research} v0.4.0 is a machine-assisted
methodology scaffold that forces these steps to happen, on the record,
before the first draft exists. The bundle does not emit filing-ready
prose. It imposes a six-step sequence (see \cref{fig:six-mode})
over the attorney's own drafting work, with hash-locked
pre-registration of search scope and a structured peer-critique pass
(severity-labeled BLOCK / WARN / NOTE, illustrated in
\cref{fig:critique}) before the draft ships.

\textbf{The paper argues} that the bundle (1) makes the attorney's
prior-art search scope reproducible by a second attorney reading the
audit trail, (2) surfaces \S112 enablement gaps as BLOCK-severity
findings before the final review pass, and (3) does so without
narrowing the claim scope the attorney would have chosen manually.
Each claim is grounded in three reproducible walkthroughs and a
four-class competitive analysis.

\textbf{The paper does not argue} that the bundle saves wall-clock
time against any baseline; no measured time-cost data exists for
v0.4.0. It does not claim to replace an attorney's legal judgment, to
improve patentability of a weak invention, or to reduce malpractice
exposure. The bundle's default hosted-LLM configuration introduces
trade-secret risks that make it inappropriate for confidential-invention
matters until an on-prem deployment mode exists that an attorney has
validated against their firm's confidentiality obligations.

\textbf{Headline limitation.} v0.4.0 has passed programmatic
verification (five tests) and simulated-persona acceptance (three
personas, all performed by the author). No real patent attorney has
used the bundle on a real filing. Read this paper as a pre-validation
capability description soliciting exactly that independent evaluation.

\textbf{Action for the reader.} Prosecution attorneys curious about
methodology-visibility tooling should run the bundle on a
non-confidential disclosure and send feedback. Firms with confidential
matter should wait for the on-prem deployment mode (roadmap: v1.x).

\begin{figure}[h]
\centering
\includegraphics[width=\linewidth]{fig-six-mode-pipeline.pdf}
\caption{The six-mode pipeline of \texttt{amplifier-bundle-research}
v0.4.0. Each mode writes a named artifact to disk before the next
mode reads it. The artifact chain (\cref{fig:artifact-chain}) is
the audit trail.}
\label{fig:six-mode}
\end{figure}

% ============================================================================
\section{Problem Statement}
% ============================================================================

\subsection{The craft attorneys already practice}

Patent attorneys drafting for USPTO examination work against three
statutory lenses that are, in substance, defensibility questions.
\cref{tab:statutes} lays out the correspondence.

\begin{table}[h]
\centering
\small
\begin{tabularx}{\linewidth}{@{}l X X@{}}
\toprule
\textbf{Statutory lens} & \textbf{The question underneath it} & \textbf{Scientific-rigor analogue} \\
\midrule
\S102 (novelty)       & Has anyone already done this?                                   & Literature review; prior-art adequacy \\
\S103 (non-obviousness)& Would the obvious next step for a person of ordinary skill get here? & Comparison-to-baseline; incremental-contribution argument \\
\S112 (enablement + written description) & Does the specification teach someone skilled how to build it? & Methods reproducibility; reader-can-rebuild test \\
\bottomrule
\end{tabularx}
\caption{The three statutory lenses patent attorneys apply, with the
scientific-rigor discipline each lens overlaps.}
\label{tab:statutes}
\end{table}

The correspondence is not accidental. Both domains evolved to protect
downstream readers (examiners, reviewers, competitors, courts) from
claims the author cannot back up. The craft is continuous even if the
vocabulary differs.

\subsection{How general-purpose AI leaks the craft}

General-purpose large language models are now commonplace in
first-draft patent work~\citep{ipwatchdog2024,patentlyo2024}. The
visible failure mode (hallucinated references) gets caught by any
attorney who reads the output. The \emph{invisible} failure mode is
the skipping of disciplined steps the human was silently carrying:

\begin{enumerate}[leftmargin=*,itemsep=0.3em]
\item \textbf{Prior-art search scope drift.} An attorney asks the
  model to draft, receives a draft, notices a reference they had not
  planned to exclude, and quietly excludes it. No audit trail records
  that the scope narrowing happened.
\item \textbf{Silent enablement gaps.} The model writes a plausible
  detailed description using intuitive-sounding parameters. The
  specification looks complete; it is not. A person of ordinary skill
  in the art (a ``PHOSITA'' in \S112 vocabulary) could not rebuild
  the invention from the description because certain assumptions were
  never stated.
\item \textbf{Adverb creep.} ``Substantially,'' ``dramatically,''
  ``significantly'' accumulate in the detailed description. Each is
  a small overclaim; together they weaken the enablement the
  specification is supposed to provide.
\end{enumerate}

None of these failures are specific to AI; they have existed as long
as patent drafting has. AI accelerates draft production without
accelerating the disciplined checks that would catch them. The
checks get deferred to the final review pass, where they cost more to
address.

\subsection{Regulatory context: USPTO AI-use disclosure}

The USPTO's 2024 AI-assistance guidance, updated in 2025, imposes
registration and disclosure obligations when AI tools are used in
prosecution, and extends the duty of candor (37 CFR 1.56) to cover
AI-induced errors~\citep{uspto2024ai}. An attorney using AI-assisted
drafting is obligated to understand (and, if necessary, disclose)
which steps the tool performed. A methodology scaffold that makes
those steps visible is adjacent to a compliance benefit. The paper
does not claim that the bundle satisfies USPTO AI-disclosure
requirements on its own; it claims only that the methodology trail
the bundle produces is structured such that disclosure is easier to
draft.

% ============================================================================
\section{Background and Market Context}
% ============================================================================

Tools that touch the ``draft patent applications'' job split into
four categories. The bundle does not replace any of these; it
occupies a niche they collectively under-serve.

\textbf{Class 1: IP portfolio and docket management suites.} PatSnap,
Anaqua, Clarivate IP Management. Mature prior-art search across
millions of records, portfolio analytics, deadline docketing, family
management. Not a drafting tool. These are comparison points on the
``auditable methodology trail'' dimension only.

\textbf{Class 2: AI-assisted drafting tools.} Rowan Patents, Solve
Intelligence, ClaimMaster, Patent Bots, and similar 2024--2025
entrants. Optimized for filing-ready prose output and claim-chart
automation. Methodology layer is proprietary and not externally
auditable. An attorney cannot reproduce a drafting decision without
vendor cooperation.

\textbf{Class 3: General-purpose LLMs.} Zero setup, universal access,
no patent-specific scaffolding. The attorney carries all discipline
manually. The de facto baseline for solo practitioners and small
firms that cannot justify Class 1 procurement.

\textbf{Class 4: Firm-internal checklists and templates.} Tacit
knowledge encoded in drafting templates and MPEP-derived checklists.
Battle-tested but invisible to audit. Scales poorly across attorneys;
degrades with turnover.

\subsection{Related prior work}

The conceptual ancestor of the bundle's approach is the
pre-registration movement in empirical science~\citep{nosek2018}. In
fields where authors commit their analysis plan to a public
repository before seeing the data, downstream rates of
motivated-reasoning post-hoc analysis drop measurably.
ClinicalTrials.gov registration obligations under the FDA Amendments
Act of 2007~\citep{fdaaa2007} extended the principle into a
regulatory framework for medical trials.

This is an analogical argument, not a transferred evidence base.
Nothing in the patent-specific literature establishes that
pre-registration discipline improves patent prosecution outcomes.
The bundle takes the \emph{mechanism} (lock methodology decisions
before results, create an audit trail for any later narrowing) and
applies it to a domain where the same failure modes are documented
and familiar.

\subsection{Development timeline}

The bundle's public history is short. v0.2.0 was released April 22,
2026 (initial public release). v0.4.0 was released the same day,
adding programmatic verification, the LaTeX compilation pipeline,
and simulated-persona acceptance tests. This paper is dated the day
after v0.4.0 and is the first public capability description. A
same-day capability paper after a tool release warrants skepticism
about commercial motivation; the primary goal of the paper is to
solicit independent evaluation, not to drive sales (the bundle is
free).

% ============================================================================
\section{Proposed Approach}
% ============================================================================

\subsection{Core concept}

The bundle imposes an auditable six-step drafting sequence over the
attorney's own work. It does not take over drafting. It provides six
sequential modes, each of which produces a concrete artifact that an
attorney (or a second attorney) can inspect, reproduce, or challenge.

\Cref{fig:six-mode} shows the mode sequence with the artifact each
mode commits to disk. \Cref{fig:agent-composition} shows the
ten specialist agents the bundle ships and which mode invokes each.
The research-coordinator agent sits above the pipeline, routing
natural-language requests to the right mode and enforcing the
honest-pivot behavior at every mode boundary.

\begin{figure}[h]
\centering
\includegraphics[width=0.95\linewidth]{fig-agent-composition.pdf}
\caption{The ten specialist agents of the research bundle, grouped by
the mode that invokes each. research-coordinator handles routing and
honest-pivot enforcement across mode boundaries. honest-critic loads
the \texttt{stop-slop} skill inline during /critique.}
\label{fig:agent-composition}
\end{figure}

\subsection{How it works}

\textbf{/question.} The attorney types the invention description in
plain English. The hypothesis-designer agent produces a structured
artifact naming a one-line statement of the invention, a claim frame
(what kind of claim it is and what it is not), directional
predictions with named mechanisms, disconfirmation criteria, and a
novelty note positioning the invention against prior-art classes.
The artifact can survive into the Field of Invention / Summary
sections of the specification.

\textbf{/study-plan.}\footnote{The slash command is /study-plan, not
/plan, because /plan is reserved by a companion bundle
(amplifier-bundle-modes). Prose elsewhere in this paper uses
``the plan mode'' as shorthand when the slash-command naming is not
the point.} Produces a hash-sealed \texttt{preregistration.yaml}
committing to the prior-art search plan (search terms, CPC/IPC
classification codes, date cutoff, excluded classes), the
enablement-demonstration plan, the honest-pivot clause (what happens
if the search reveals a claim needs narrowing), and abandonment
criteria.

This is the load-bearing innovation. A hash-sealed search plan means
any later narrowing (excluding a class because it surfaced
uncomfortable art, dropping a search term because it returned too
much to review) shows up in a timestamped amendment log rather than
hiding in the final draft. A second attorney auditing the work can
reproduce the planned search and compare it to the executed search.
The cost of reproducibility is paid at plan time, not at
office-action-response time.

\textbf{/execute.} Runs the prior-art search against configured
sources (USPTO, Google Patents, arXiv, IEEE Xplore), populates a
claim chart, and runs an enablement demonstration if one was planned.
Any deviation from the plan goes into the honest-pivot log. The
output is an evidence log keyed to the claims /question identified.

\textbf{/critique.} The honest-critic agent issues severity-labeled
findings across five dimensions: thesis defensibility, overclaiming
detection, evidence gaps, generalization limits, and limitation
specificity. BLOCK is reserved for issues that would survive into the
filed application as traceable weaknesses (silently narrowed claims,
unsupported enablement statements, superlatives without numeric
backing). During /critique, honest-critic loads the stop-slop
skill~\citep{pandya2024stopslop} inline to detect AI-tell writing
patterns (adverb piles, binary contrasts, false agency) that arrive
through general-LLM-generated prose. The attorney retains
accept-reject control over every finding; the bundle emits the
finding-with-context, not the forced edit. \Cref{fig:critique}
shows how findings route through the /publish gate.

\begin{figure}[h]
\centering
\includegraphics[width=0.85\linewidth]{fig-critique-severity-flow.pdf}
\caption{How /critique findings route through the /publish gate.
BLOCK findings must be resolved before the draft proceeds; WARN and
NOTE findings are surfaced in the final PDF's limitations section.}
\label{fig:critique}
\end{figure}

\textbf{/draft.} Populates the venue-specific template. For patent
work, the template is the USPTO invention-disclosure structure
(Field of Invention, Background, Summary, Detailed Description,
Claims, Abstract, Drawings). The technical-writer agent composes the
sections from the evidence log; the citation-manager agent assembles
the prior-art bibliography in USPTO format.

\textbf{/publish.} The venue-formatter agent applies USPTO document
formatting (37 CFR 1.52 margins, line spacing, consecutive page
numbering, claim numbering) and emits filing-ready DOCX and PDF. A
deterministic Python script (\texttt{scripts/compile\_latex.py})
handles the final compilation step where mechanical correctness
matters more than language-model fluency. The /publish gate refuses
to execute unless at least one /critique pass has been run and any
honest-pivot entries have been acknowledged.

\begin{figure}[h]
\centering
\includegraphics[width=\linewidth]{fig-artifact-chain.pdf}
\caption{The chain of session artifacts produced by a single
/question--to--/publish run. The sha256 seal on
\texttt{02-study-plan.yaml} is the integrity mechanism for
honest-pivot enforcement; /publish recomputes it and blocks on
mismatch.}
\label{fig:artifact-chain}
\end{figure}

\subsection{What makes it different}

Most AI-assisted drafting tools (Class 2) and general LLMs (Class 3)
optimize for \emph{output quality}. The bundle optimizes for
\emph{methodology visibility}. The distinguishing choices:

\begin{enumerate}[leftmargin=*,itemsep=0.2em]
\item Pre-registration is mandatory, not opt-in. /publish blocks if
  /study-plan did not produce a sealed prereg.
\item /critique is a separate, structured pass with labeled
  findings an attorney can re-inspect a year later.
\item The attorney holds the pen at every mode boundary. The bundle
  never commits an edit without acceptance.
\item The audit trail is human-readable YAML. A second attorney can
  reproduce a run from the session files.
\item Output is inspectable at every stage. Every mode writes a
  named artifact to disk. A failed run is debuggable from those
  artifacts.
\end{enumerate}

These choices come at a cost. The bundle runs slower than a raw-LLM
``draft me this application'' prompt because the sequence takes as
long as the attorney spends reviewing each mode's output. For an
attorney whose pain is ``I want prose fast,'' the bundle is the
wrong tool. For an attorney whose pain is ``I can write the prose; I
cannot remember every step I should have done,'' the bundle's
sequence carries the discipline.

% ============================================================================
\section{Evidence and Case Studies}
% ============================================================================

All three case studies are author-generated: two dogfood runs of the
bundle on itself, one simulated-persona acceptance test the author
ran during v0.4.0 development. They are capability demonstrations
(the bundle produces these artifacts when run), not outcome evidence
(these artifacts improved attorney filings). Any attempt to present
these as outcome evidence fails the paper's own /critique pass.

\subsection{Case study 1: /question dogfood}

\textbf{Context.} The bundle was asked to run /question on the
statement ``amplifier-bundle-research should be valuable for patent
attorneys.'' Captured during v0.3 verification on 2026-04-22 (git
commit \texttt{d528196}).

\textbf{Outcome.} A 97-line structured \texttt{sharpened\_question.yaml}
containing a one-line thesis with explicit claim type, three
directional predictions with measurable criteria, disconfirmation
conditions, scope boundaries, and a four-class novelty note.

\textbf{What this does and does not prove.} It proves the bundle
produces a question-sharpening artifact with the claimed shape. It
does \emph{not} prove the artifact is useful to a patent attorney;
no attorney has reviewed it.

\subsection{Case study 2: /study-plan smoke test}

\textbf{Context.} A /study-plan run on a synthetic software-ML
invention, producing a populated \texttt{patent-prereg.yaml}.
Captured during v0.3 verification.

\textbf{Outcome.} A 57-line prereg artifact containing the prior-art
search plan, enablement-demonstration plan, honest-pivot clause,
abandonment criteria, and sha256 seal.

\textbf{What this does and does not prove.} It proves the pipeline
produces a structurally complete pre-registration artifact. It does
\emph{not} prove the search plan is well-chosen for any real
invention, nor that an attorney would accept the generated terms
without revision.

\subsection{Case study 3: simulated Persona-A acceptance test
(the mixed case)}

\textbf{Context.} The v0.4 simulated-persona acceptance test. The
author played the role of a patent attorney taking a plausible
software-ML invention through /question and /study-plan.

\textbf{Outcome.} A patent-schema preregistration with prior-art
search plan, enablement demonstration, honest-pivot clause, and
abandonment criteria. Overclaiming adverbs in the scratch
specification were flagged for removal during the pass.

\textbf{What this does and does not prove (mixed case).} The author
is both the bundle designer and the simulated attorney. This is a
correctness-of-shape test, not a usefulness test. A real attorney
reading the output might find it overengineered, underengineered, or
calibrated wrongly in ways the author cannot self-detect. The case
is cited as evidence of capability, not of usefulness. The gap
between capability and usefulness is the paper's headline
limitation, not a footnote.

\subsection{Dogfood of this paper itself}

The paper you are reading was produced by running the white-paper
recipe over the /question--/publish pipeline. Every mode wrote an
artifact to \texttt{output/whitepaper-run/}
(\cref{fig:artifact-chain}). The /critique pass produced fourteen
findings: four BLOCK, six WARN, four NOTE, plus four named
alternative explanations. All BLOCK resolutions are visible in this
draft (see especially the first-person disclosure in the front
matter and the headline-limitation framing in the Executive
Summary). The second /critique pass, run after loading the
\texttt{stop-slop} skill, flagged seven additional prose-level
patterns; all were resolved before compilation.

% ============================================================================
\section{Competitive Analysis}
% ============================================================================

The bundle occupies a methodology-visibility niche.
\cref{tab:competitive} compares along a single axis: does the
workflow produce an attorney-reviewable audit trail of its
methodology decisions? Classes 1 and 4 are adjacent-concern rather
than rival products. They appear here for completeness, not
to suggest the bundle outperforms them on their native axes.

\begin{table}[h]
\centering
\scriptsize
\setlength\tabcolsep{3pt}
\begin{tabularx}{\linewidth}{@{}p{2.1cm} X X p{2.4cm}@{}}
\toprule
\textbf{Approach} & \textbf{Strengths (methodology-visibility niche)} & \textbf{Weaknesses (same niche)} & \textbf{Best fit} \\
\midrule
IP management suites (PatSnap, Anaqua, Clarivate)
  & Mature prior-art search; portfolio and docket management; enterprise audit-readiness
  & Not a drafting tool; no pre-registration of search scope; priced for enterprise
  & Firms with portfolio-scale docketing and procurement budget \\
\midrule
AI-assisted drafting (Rowan, Solve Intelligence, ClaimMaster, Patent Bots)
  & Filing-ready prose output; claim-chart automation; integrated UI
  & Methodology layer proprietary and not externally auditable; attorney cannot reproduce drafting decisions
  & Attorneys whose primary pain is prose-output speed \\
\midrule
General-purpose LLMs (ChatGPT, Claude, Gemini)
  & Zero setup; universal access; wide general capability
  & Attorney carries all discipline manually; no audit trail; inconsistent across sessions
  & Quick drafting tasks where the attorney provides the methodology layer \\
\midrule
Firm-internal checklists / AIPLA templates
  & Battle-tested; tuned to firm philosophy; free
  & Invisible to audit; scales poorly; degrades with turnover; rarely portable
  & Firms with strong senior-to-junior tacit transmission \\
\midrule
\textbf{\texttt{amplifier-bundle-research} v0.4.0 (this paper)}
  & Pre-registration of search scope before execution; severity-labeled critique pass; honest-pivot enforcement; USPTO-compliant venue formatting; open source; one-command install
  & \textbf{No real-attorney UX validation yet} (v0.4 is simulated-persona only); no measured time-cost comparisons; trade-secret-sensitive matters unsafe with default hosted-LLM configuration; prose-polish less optimized than Class-2 competitors
  & Solo / small-firm attorneys who want methodology visibility and accept the trade of prose speed for audit trail \\
\bottomrule
\end{tabularx}
\caption{Competitive comparison along a single axis: produces an
attorney-reviewable audit trail of methodology decisions. The
bundle's honest weaknesses row is populated deliberately.}
\label{tab:competitive}
\end{table}

\subsection{Alternative explanations a skeptical reviewer would raise}

\begin{enumerate}[leftmargin=*,itemsep=0.2em]
\item \textbf{Novelty effect.} Attorneys may find the bundle appealing
  because it is new. Longitudinal real-attorney use would distinguish
  this from durable appeal.
\item \textbf{Well-written-checklist effect.} The methodology-visibility
  benefit may be reachable by a sufficiently well-written paper
  checklist. A three-arm comparison (raw LLM / paper checklist /
  bundle) would be needed to isolate the bundle's contribution. That
  comparison has not been run.
\item \textbf{Boilerplate acceleration confound.} Any observed time
  savings may be attributable to LLM acceleration of boilerplate
  prose (a capability all AI-drafting tools share) rather than to
  the bundle's methodology scaffolding.
\item \textbf{Over-narrowing risk.} /critique could encourage
  attorneys to over-narrow claims by surfacing every possible
  concern. Prediction P3 (claim-scope preservation) is an attempt to
  forestall this; it is not yet empirically validated.
\end{enumerate}

% ============================================================================
\section{Implementation Considerations}
% ============================================================================

\textbf{Prerequisites.} An installation of Amplifier (open source,
one-line install); access to a language-model provider configured
through Amplifier's provider layer (Anthropic Claude, OpenAI, Azure
OpenAI, or local Ollama); familiarity with the six mode names and
order. Plain-language explanations appear in the /study-plan mode's
own output when the attorney runs it.

\textbf{Deployment path.} A realistic adoption arc for a solo
practitioner: (1) pilot on a non-confidential disclosure, (2) inspect
the audit trail files, (3) assess critique-finding quality, (4) do
not move to confidential matter until the on-prem deployment mode is
validated.

\textbf{Integration points.} The bundle operates on text files on
the attorney's local filesystem. It does not require changes to the
firm's docketing system, PDM, or document-management workflow.
Output is DOCX and PDF in standard USPTO format.

\textbf{Cost and effort.} The bundle is free (Apache-2.0). The
meaningful cost is attorney time to learn the six-mode vocabulary
and review each artifact. Order-of-magnitude estimate: half a day of
hands-on learning on a pilot disclosure before an attorney can judge
fit. Real-attorney validation has not measured this estimate. The
paper does not claim the bundle saves attorney time compared to any
baseline.

% ============================================================================
\section{Risk Analysis}
% ============================================================================

\begin{table}[h]
\centering
\footnotesize
\begin{tabularx}{\linewidth}{@{}p{3.1cm} p{1.6cm} p{1.6cm} X@{}}
\toprule
\textbf{Risk} & \textbf{Likelihood} & \textbf{Impact} & \textbf{Residual (after mitigation)} \\
\midrule
\textbf{R1. Trade-secret leakage.} Default configuration sends invention disclosures to a hosted provider; may violate client confidentiality or (rarely) export-control restrictions.
  & Medium & \textbf{High} & \textbf{Significant. Unresolved for v0.4.0.} Do not use default configuration on trade-secret matter until on-prem deployment is validated (v1.x roadmap). \\
\midrule
\textbf{R2. Attorney-UX unvalidated.} v0.4.0 has passed simulated-persona acceptance only. No real patent attorney has used it on a real filing.
  & High & Medium & \textbf{Significant. The headline limitation.} Early adopters should treat their own use as pilot data. \\
\midrule
\textbf{R3. Claim-scope regression.} honest-critic may flag
  defensible language as overclaiming, leading an attorney under
  deadline pressure to accept narrower claims than they should file.
  & Medium & Medium & Minor. Mitigable with normal review practice; v1.0 validation will test accept/reject behavior under realistic deadline conditions. \\
\midrule
\textbf{R4. Methodology theater.} The bundle could become ceremony of rigor rather than actual rigor, providing false assurance.
  & Medium & Medium--High & Minor but requires firm-level culture of honest use. Audit-trail design makes skipping visible after the fact. \\
\bottomrule
\end{tabularx}
\caption{Risk analysis. R1 and R2 remain significant for v0.4.0. The
honest framing: v0.4.0 is appropriate for non-confidential pilot use
by attorneys willing to treat their own usage as validation data.}
\label{tab:risks}
\end{table}

% ============================================================================
\section{Limitations}
% ============================================================================

\begin{description}[leftmargin=1.5em,itemsep=0.3em]
\item[L1. No real-attorney validation yet.] v0.4.0 passed programmatic
  verification (five tests) and simulated-persona acceptance (three
  personas, all performed by the author). No real patent attorney
  has used the bundle on a real filing. All capability claims in
  this paper are of the form ``the bundle produces this artifact
  when run,'' not ``this artifact improved an attorney's filing.''

\item[L2. No measured time-cost comparisons.] The paper makes no
  claim that the bundle is faster or slower than any baseline. No
  wall-clock data exists on either side of any comparison.

\item[L3. Default configuration inappropriate for trade-secret matter.]
  See R1. You should not use the default hosted-LLM configuration on
  trade-secret-sensitive matters until an on-prem deployment mode
  exists that you have validated against your firm's confidentiality
  obligations.

\item[L4. Analogical-reasoning load.] The pre-registration mechanism
  draws from an evidence base in empirical science~\citep{nosek2018}
  that does not directly apply to patent prosecution. The argument
  is mechanism-similarity, not transferred evidence.

\item[L5. Competitive analysis is asymmetric by design.] Classes 1
  and 4 are adjacent product categories, not direct rivals. The
  table compares on a single axis (methodology-visibility audit
  trail). Readers choosing among these products on other axes
  (portfolio management, prose polish, firm-tooling integration)
  should not use \cref{tab:competitive} as a ranking.

\item[L6. The author wrote the paper about the author's tool.] Every
  capability claim should be read with this conflict in mind. The
  paper is an invitation to independent evaluation, not an
  independent evaluation.

\item[L7. Vocabulary drift.] ``Pre-registration,'' ``hash-locked
  plan,'' and ``sealed methodology'' are used across bundle
  documentation to refer to the same mechanism. This paper uses
  ``pre-registration'' consistently.
\end{description}

% ============================================================================
\section{Recommendations}
% ============================================================================

\textbf{For solo practitioners and small firms.} Run the bundle on a
non-confidential disclosure. Open the artifact files. Decide whether
the methodology scaffolding matches your own drafting discipline.
Feed back what does not fit. Do not use the default configuration on
confidential matter until the on-prem deployment mode is validated.

\textbf{For firms with confidential matter and procurement budget.}
Treat v0.4.0 as a reference architecture for methodology visibility.
The open-source nature of the bundle allows the methodology layer to
be re-implemented behind a firm-internal deployment. Procurement-scale
decisions should wait for the v1.x on-prem path and for
real-attorney validation data.

\textbf{For attorneys considering participating in v1.0 validation.}
The path is cheap: a one-line install and a non-confidential
disclosure. Participation would produce exactly the evidence this
paper cannot yet provide.

% ============================================================================
\section{Conclusion}
% ============================================================================

\texttt{amplifier-bundle-research} v0.4.0 imposes a six-step,
auditable sequence over a patent attorney's first-draft work. It
makes the attorney's own methodology visible, reproducible, and
portable. It does not replace legal judgment, does not emit
filing-ready prose autonomously, and \textbf{has not yet been
validated by a real attorney on a real filing}.

The strongest evidence is three reproducible walkthroughs (CS1, CS2,
CS3) that demonstrate the bundle produces artifacts of the claimed
shape, and five programmatic verifications that document the tool
protocol, compile pipeline, and dogfood paths work end-to-end. The
strongest honest concession is that every capability claim is a
capability claim, not an outcome claim. No data exists on whether the
bundle helps attorneys write better applications in fewer cycles.
That gap is the paper's reason for existing: to solicit the
real-attorney sessions that would produce that data.

Attorneys curious about methodology-visibility tooling should install
the bundle, run it on a non-confidential disclosure, and send
feedback. The repository is public, the source is readable, and the
v1.0 validation milestone is open for participation.

% ============================================================================
\section*{About the Author}
% ============================================================================

Michael Jabbour authored \texttt{amplifier-bundle-research}. He is
not a practicing patent attorney. His prior work includes the design
of methodology-enforcement tooling for evidence-defensibility
domains (scientific rigor, technical due diligence). The bundle
described in this paper is an attempt to carry methodology-enforcement
into patent-drafting workflows. He benefits reputationally (not
financially in v0.4.0) from adoption of the bundle; the paper is a
first-person capability description soliciting independent
evaluation.

% ============================================================================
\bibliographystyle{plainnat}
\begin{thebibliography}{11}

\bibitem[IPWatchdog(2024)]{ipwatchdog2024}
IPWatchdog.
\newblock AI-Assisted Patent Drafting: Practitioner Perspectives,
  2024--2025.
\newblock Series of articles archived at \url{https://ipwatchdog.com/}
  (search: ``AI drafting 2024'').

\bibitem[Patently-O(2024)]{patentlyo2024}
Patently-O.
\newblock Artificial Intelligence and Patent Practice.
\newblock Series of posts by Dennis Crouch et al., 2024--2025,
  archived at \url{https://patentlyo.com/}.

\bibitem[USPTO(2024)]{uspto2024ai}
United States Patent and Trademark Office.
\newblock \emph{Guidance on Use of Artificial Intelligence-Based Tools
  in Practice Before the USPTO}.
\newblock April 2024, with subsequent 2025 updates.
\newblock \url{https://www.uspto.gov/} (search: ``AI guidance 2024'').

\bibitem[Nosek et~al.(2018)]{nosek2018}
B.~A. Nosek, C.~R. Ebersole, A.~C. DeHaven, and D.~T. Mellor.
\newblock The preregistration revolution.
\newblock \emph{Proceedings of the National Academy of Sciences},
  115(11):2600--2606, 2018.
\newblock \url{https://doi.org/10.1073/pnas.1708274114}.
\newblock \textbf{Cited as analogical support only.}

\bibitem[USPTO-MPEP(2024)]{uspto_mpep}
United States Patent and Trademark Office.
\newblock \emph{Manual of Patent Examining Procedure (MPEP)},
  \S2100 series (patentability), \S2141 (obviousness / Graham
  factors), \S2163 (written description and enablement under
  \S112(a)).
\newblock \url{https://www.uspto.gov/web/offices/pac/mpep/}.

\bibitem[37 CFR 1.52]{cfr152}
37 CFR 1.52.
\newblock Language, paper, writing, margins, compact disc
  specifications.
\newblock Governs USPTO document-formatting requirements referenced
  in the venue-formatter implementation.

\bibitem[FDAAA(2007)]{fdaaa2007}
FDA Amendments Act of 2007, Public Law 110-85, \S801.
\newblock Establishes the ClinicalTrials.gov registration obligation
  cited as an analogical regulatory-pre-registration precedent.

\bibitem[Graham(1966)]{graham1966}
Graham v. John Deere Co.
\newblock 383 U.S. 1 (1966).
\newblock Supreme Court decision establishing the four-factor
  obviousness test referenced under \S103.

\bibitem[Jabbour(2026)]{jabbour2026repo}
M.~Jabbour.
\newblock \emph{amplifier-bundle-research} repository and release
  notes.
\newblock \url{https://github.com/michaeljabbour/amplifier-bundle-research}.
\newblock v0.2.0 (commit \texttt{fc8ba42}, Apr 22 2026); v0.4.0
  (commit \texttt{d528196}, Apr 22 2026).

\bibitem[K-Dense(2025)]{kdense2025}
K-Dense AI.
\newblock \emph{scientific-agent-skills} repository.
\newblock \url{https://github.com/K-Dense-AI/scientific-agent-skills}.
\newblock Upstream source of the pre-registration, peer-review, and
  scholar-evaluation skill primitives wrapped by the bundle's agents.

\bibitem[Pandya(2024)]{pandya2024stopslop}
H.~Pandya.
\newblock \emph{stop-slop}: a skill file for removing AI tells from
  prose.
\newblock \url{https://github.com/hardikpandya/stop-slop}.
\newblock Adapted to Amplifier form for this bundle; see
  \texttt{.amplifier/skills/stop-slop/}.

\end{thebibliography}

\end{document}
